\documentclass[11pt]{article}
\usepackage{amsmath,amssymb}
\usepackage[margin=1in]{geometry}
\usepackage{booktabs}
\usepackage{tabularx}
\usepackage{hyperref}
\setlength{\emergencystretch}{2em}

\title{General Open-World Scientific Discovery Superiority:\\
Discover More, Preserve Search Optionality, and Know When the Science Is Closed}
\author{ORION Paper II Ultimate-Successor Working Manuscript}
\date{August 2026}

\begin{document}
\maketitle

\begin{abstract}
Deep-research agents increasingly combine iterative retrieval, planning, citation, coverage estimation, state management, and evidence-aware termination. ORION Paper II therefore targets the strongest upward claim suggested by this ecosystem: across heterogeneous scientific domains, route structures, evidence regimes, and changing questions, can ORION outperform a donor-complete deep-research system at discovering decision-relevant evidence, preserving future search optionality, synthesizing contradictions, and issuing globally calibrated scientific closure under matched total resources? Existing P2 adverse and provider-invalid probes are treated as causal maps of where open-world research can fail, not as reasons to retreat to a local stopping-rule claim. The successor separates route generation, retrieval, inclusion, processing, state representation, optionality, provider validity, and task-global closure, while allowing the strongest comparator to contain modern retrieval, evidence-aware stopping, coverage estimation, contradiction handling, caching, and query-conditioned state compilation. Protected tasks hide distributed evidence and future route changes, so success requires both finding what matters now and retaining enough state to discover what matters later. Primary outcomes include decision-relevant recall and precision, final scientific utility, false closure, clean closure, unresolved calibration, optionality loss, cross-domain transfer, and total cost. The intended terminal is replicated open-world discovery and closure superiority over the strongest donor-complete system, ideally including at least one new exploration or closure mechanic discovered from a prior loss/tie that expands protected reach. Existing negative and \textsc{CannotCheck} results remain immutable and drive upward successor search rather than default claim contraction.
\end{abstract}

\section{Largest scientific question}
The target is not a narrow residual after modern deep-research systems. It is:
\begin{quote}
At matched model, tool, information, and total resource budgets, can ORION become a generally superior open-world scientific discovery process by finding more of the evidence that matters, preserving the ability to discover later-relevant evidence, and closing only when the task-level scientific obligations are actually discharged?
\end{quote}

DeepResearch Bench evaluates research-report quality and retrieval/citation behavior on expert-authored tasks \cite{deepresearchbench2026}. Evidence-aware termination has been proposed explicitly to reduce premature stopping \cite{dontstopearly2026}. SAGE shows that reasoning-intensive scientific retrieval remains difficult and can invert expectations about which retrievers are strongest \cite{sage2026}. These works raise the baseline. P2-U absorbs their mechanisms into the donor-complete comparator and moves upward to a stronger question: whether ORION can coordinate retrieval, processing, state construction, exploration, and closure better across regimes than any fixed composition of those donor mechanics.

\section{Upward-generalization principle}
Any donor mechanism that solves a P2-U subclass is added to the comparator. ORION earns superiority only at the broader level, through one or more of:
\begin{enumerate}
    \item higher decision-relevant discovery under matched total resources;
    \item lower false global closure at equal or higher scientific utility;
    \item better adaptation to heterogeneous route/evidence regimes;
    \item preserved optionality under changing questions and evidence horizons;
    \item a learned or invented exploration/closure mechanic that expands protected reach beyond the frozen donor-complete basis;
    \item cross-domain transfer of the same governing search-and-closure policy.
\end{enumerate}
The intended terminal is \textbf{GENERAL_OPEN_WORLD_DISCOVERY_SUPERIORITY}. Narrower stage-local results are fallback scientific outcomes, not the default research objective.

\section{Immutable adverse evidence}
The predecessor P2 external probes include non-positive outcomes and provider-invalid cases. They are retained as stage-localizing evidence, not rewritten as performance wins or losses. The successor preserves these distinctions:
\begin{enumerate}
    \item evidence never retrieved;
    \item evidence retrieved but excluded incorrectly;
    \item evidence included but not processed into a decision-usable state;
    \item a material route never generated or revisited;
    \item provider/tool output invalid or unavailable;
    \item local route closure incorrectly promoted to task-global closure.
\end{enumerate}
A provider-invalid run is \textsc{CannotCheck}. It is never scored as a comparator failure. These failures define dimensions over which the larger search policy must generalize.

\section{Task model}
Each protected task contains a hidden evidence universe partitioned into material routes $r\in R$. A route may contain supportive, contradictory, redundant, superseded, or non-decision-relevant evidence. Candidate systems do not receive $R$ or the gold materiality labels.

At time $t$, the system maintains:
\[
S_t=(Q_t,E_t,P_t,U_t,C_t),
\]
where $Q_t$ is the active query/route frontier, $E_t$ is acquired raw evidence, $P_t$ is processed evidence state, $U_t$ is unresolved/optionality state, and $C_t$ is the candidate closure status. The protected evaluator separately knows whether material evidence or routes remain open.

\section{Closure states}
The only task-level terminals are:
\begin{itemize}
    \item \textsc{Closed}: all frozen closure obligations are discharged;
    \item \textsc{Open}: at least one identified material obligation remains actionable;
    \item \textsc{Unresolved}: available evidence cannot establish closure or a specific remaining route under the budget;
    \item \textsc{CannotCheck}: required provider, evidence-custody, or evaluator authority is unavailable.
\end{itemize}
\textsc{Unresolved} is a calibrated epistemic state. \textsc{CannotCheck} is an authority/evidence-availability state. Neither is silently converted into \textsc{Closed}.

\section{Prospective benchmark families}
The benchmark is frozen before protected outcomes and includes paired constructions that isolate the following failures.

\begin{table}[h]
\centering
\small
\begin{tabularx}{\textwidth}{@{}lXX@{}}
\toprule
Family & Hidden condition & Failure exposed \\
\midrule
Missing route & decisive source exists behind an ungenerated route & route generation/coverage \\
Retrieved but unprocessed & decisive source was fetched but not incorporated & processing state \\
Duplicate evidence & many citations map to one underlying evidence unit & false coverage inflation \\
Contradictory sources & high-quality sources disagree materially & synthesis/closure calibration \\
Provider invalid & mandatory route returns invalid/unverifiable output & authority, not performance \\
Apparently saturated & repeated retrieval yields duplicates while hidden route remains & premature stopping \\
Optionality loss & compressed state answers current query but erases future route cue & state preservation \\
Clean closure & all obligations truly discharged & excessive abstention penalty \\
\bottomrule
\end{tabularx}
\caption{Mandatory P2-U stage-identifying families.}
\end{table}

The final study must scale beyond these authored mechanisms to naturalistic scientific tasks spanning multiple domains, source ecosystems, query shifts, contradiction structures, and at least one systematic-review-like setting with independently curated evidence relevance. At least one evaluation block changes both domain and evidence-route geometry.

\section{State optionality as a charged resource}
P2-U compares four evidence-state policies under matched total cost:
\begin{enumerate}
    \item raw accumulated history;
    \item universal evidence index;
    \item query-conditioned compiled evidence state;
    \item compile+cache/materialize policy that may preserve selected raw evidence for future routes.
\end{enumerate}
Compilation is never free preprocessing. Its token, compute, storage, tool, and latency costs are charged to the same budget as retrieval and reasoning.

Future query/route changes are frozen before the current state is built. This permits a direct optionality test: a state can preserve current answer quality yet be harmful if it deletes the only coordinate needed to detect a later material route.

\section{Baselines}
Comparators receive equivalent model/tool access and total resource budgets. The mandatory set includes:
\begin{enumerate}
    \item fixed-budget search+RAG;
    \item iterative deep-research agent;
    \item evidence-aware termination system in the style of modern EDR \cite{dontstopearly2026};
    \item strong retrieval-focused deep-research system evaluated on scientific retrieval, informed by SAGE \cite{sage2026};
    \item coverage estimator / capture-recapture augmented system;
    \item contradiction-aware evidence synthesis;
    \item query-conditioned state compilation/cache management;
    \item donor-complete product containing all preceding mechanisms and allowed to route among them;
    \item strongest current deep-research/scientific-search agent available at freeze time;
    \item an oracle route ceiling used only for analysis.
\end{enumerate}
Official unavailable systems are not fabricated. Mechanism-matched reimplementations are labelled as such.

\section{Primary hypotheses}
\paragraph{H1, generalized discovery superiority.}
Across the heterogeneous benchmark, ORION improves protected final scientific utility and decision-relevant evidence recall over the strongest runnable donor-complete comparator by prospectively frozen margins.

\paragraph{H2, closure superiority.}
ORION reduces false task-global closure while preserving non-inferior or superior clean closure.

\paragraph{H3, optionality superiority.}
ORION preserves future evidence/search optionality better under changing questions at matched total cost.

\paragraph{H4, cross-regime transfer.}
The H1--H3 advantage survives held-out domains, source ecosystems, route geometries, and query shifts under independent route accounting.

\paragraph{H5, mechanism expansion.}
At least one negative/tied family yields a new exploration, state, or closure mechanic that expands protected reach beyond the pre-freeze donor-complete basis and transfers to fresh tasks.

\section{Outcomes and statistical plan}
Every task is an independent experimental unit unless a hierarchy is explicitly frozen. Repeated model samples within one task are technical repeats and are not counted as independent $n$.

Primary outcomes are:
\begin{enumerate}
    \item decision-relevant evidence recall and precision;
    \item protected final answer/scientific utility;
    \item false closure rate;
    \item true clean-closure coverage;
    \item calibrated \textsc{Unresolved} decisions;
    \item route-level missed-obligation count;
    \item optionality loss under frozen future queries;
    \item cross-domain/source/route transfer;
    \item tokens, retrieval calls, state-compilation cost, wall time, and time-to-valid-closure.
\end{enumerate}

Before protected access, freeze superiority/non-inferiority margins, family weights, confidence-interval method, multiplicity policy, missing/provider-invalid handling, and catastrophic false-closure rules. Report effect sizes and uncertainty, not significance language alone.

\section{Route and task closure receipts}
A route receipt records route identity, queries attempted, evidence identities, processing status, unresolved obligations, and stop reason. A task receipt records which route obligations support global closure and which remain open. A route may be locally saturated without authorizing task closure.

A source citation count, context length, number of search calls, or repeated duplicate retrieval never constitutes a closure proof.

\section{Hostile controls}
Mandatory attacks include:
\begin{enumerate}
    \item a high citation count generated entirely from redundant evidence;
    \item a decisive document retrieved early but omitted from processed state;
    \item a late contradictory source reachable only through a distinct route;
    \item state compression that preserves the current answer but deletes a future-route cue;
    \item an unavailable provider that would otherwise contain material evidence;
    \item a clean closed case to defeat deny-all/never-stop policies;
    \item evidence-label or route-name leakage into candidate-visible metadata;
    \item a deeper-search arm showing whether extra retrieval alone closes the apparent ORION advantage;
    \item a donor-complete arm that dynamically selects among retrieval, stopping, contradiction, and state-management mechanisms.
\end{enumerate}

\section{Negative-result recursion as upward search}
Every false closure, missed material source, harmful state compilation, and tied/negative family is retained. Each is assigned one primary causal responsibility from
\[
\{\text{route generation},\text{retrieval},\text{inclusion},\text{processing},\text{state optionality},\text{stopping},\text{provider},\text{global authority},\text{benchmark}\}.
\]
The recovery programme does not stop at one-factor repair. It first freezes a discriminator to localize the failure, then searches materially different successor mechanisms and parent disciplines for a rule that generalizes beyond the originating case. A retrieval miss can trigger a new route-generation family; a state-optionality loss can trigger a compile/cache/recover policy; recurrent failures across both may justify a higher-order budget-allocation or search-state mechanic. Provider failure remains \textsc{CannotCheck} unless authority is restored.

ORION issue \#650 is the authority-bearing P2-U research programme. A child issue is warranted only for a newly observed causal stage or materially new mechanism family. Outcome deletion, margin relaxation, baseline weakening, or scoring provider failures as comparator losses are prohibited.

\section{Claim ladder}
\begin{enumerate}
    \item \textbf{Protocol frozen}: no performance claim.
    \item \textbf{Protected open-world superiority}: H1 and H2 pass with clean-closure safety.
    \item \textbf{Cross-regime discovery superiority}: H1--H4 pass across held-out domains/source/route regimes.
    \item \textbf{General open-world discovery superiority}: cross-regime superiority plus H5 demonstrates expansion of the search/closure mechanism basis from failures.
\end{enumerate}
A narrower terminal is an honest fallback only if the largest claim is not earned.

\section{Current status}
This is a prospective TeX successor. It does not modify predecessor P2 results and does not assert generalized open-world superiority before the frozen protected campaign is executed.

\begin{thebibliography}{9}
\bibitem{deepresearchbench2026}
Y. Yang et al., ``DeepResearch Bench: A Comprehensive Benchmark for Deep Research Agents,'' ICLR 2026.

\bibitem{dontstopearly2026}
P. K. Choubey et al., ``Don't Stop Early: Scalable Enterprise Deep Research with Controlled Information Flow and Evidence-Aware Termination,'' Proceedings of ACL 2026, Industry Track, pp. 1699--1713, doi:10.18653/v1/2026.acl-industry.116.

\bibitem{sage2026}
T. Hu, Y. Zhao, C. Zhang, A. Cohan, and C. Zhao, ``SAGE: Benchmarking and Improving Retrieval for Deep Research Agents,'' arXiv:2602.05975, 2026.
\end{thebibliography}

\end{document}
